% Beginning of `inweb-macros.tex`
%
% These are TeX macros for use by Inweb when rendering to a Plain TeX variant
% (such as TeX itself, pdftex or xetex): they should work on any TeX. Macros
% which do not have the luxury of engine-independence have been hived off into
% the files `links-macros.tex`, `colours-macros.tex`, and `images-macros.tex`.
%
% (1) Fonts
% There are more fonts here than we actually need for Inweb purposes, but it
% seems easiest to follow Knuth's semi-standard book design code.
%
% See Appendix B of the TeXBook: some `\font` definitions, such as `\tenrm`
% and `\tentt`, are already defined in Plain TeX.
%
\font\defnannotationfont=cmti9
\def\defnannotationtext#1{{\defnannotationfont #1}}
\font\tentex=cmtex10
\font\sourcefont=cmss10
\font\holonfont=cmss9
\def\holonnametext#1{{\holonfont #1}}
\hyphenchar\holonfont=-1
\font\usagefont=cmss8
\def\endnotetext#1{{\usagefont #1}}
\def\sectiontoctext#1{{\usagefont #1}}
\font\inchhighfont=cminch
\font\reducedinchhighfont=cminch at 60pt
\font\chaptertitlefont=cmss10 at 24pt
\def\chaptertitletext#1{\chaptertitlefont #1}
\font\sectiontitlefont=cmss10 at 20pt
\def\sectiontitletext#1{\sectiontitlefont #1}
\font\ssectiontitlefont=cmss10 at 20pt
\font\sssectiontitlefont=cmr10 at 12pt
\font\ninerm=cmr9
\font\eightrm=cmr8
\font\sixrm=cmr6
\font\ninei=cmmi9
\skewchar\ninei='177
\font\eighti=cmmi8
\skewchar\eighti='177
\font\sixi=cmmi6
\skewchar\sixi='177
\font\ninesy=cmsy9
\skewchar\ninesy='60
\font\eightsy=cmsy8
\skewchar\eightsy='60
\font\sixsy=cmsy6
\skewchar\sixsy='60
\font\holonnumberfont=cmss10 at 7pt
\def\holonnumbertext#1{{\holonnumberfont #1}}
\font\eightss=cmssq8
\font\eightssi=cmssqi8
\font\ninebf=cmbx9
\font\eightbf=cmbx8
\font\sixbf=cmbx6
\font\ninett=cmtt9
\hyphenchar\ninett=-1
\font\eighttt=cmtt8
\hyphenchar\eighttt=-1
\hyphenchar\tentt=-1 % inhibit hyphenation in typewriter type
\font\ninesl=cmsl9
\font\eightsl=cmsl8
\font\nineit=cmti9
\font\eightit=cmti8
\font\tenu=cmu10 % unslanted text italic
\font\magnifiedfiverm=cmr5 at 10pt
\font\manual=manfnt % font used for the METAFONT logo, etc.
\font\cmman=cmman % font used for miscellaneous Computer Modern variations
%
% Knuth's standard way to change point size:
%
\newskip\ttglue
\def\tenpoint{\def\rm{\fam0\tenrm}%
  \textfont0=\tenrm \scriptfont0=\sevenrm \scriptscriptfont0=\fiverm
  \textfont1=\teni \scriptfont1=\seveni \scriptscriptfont1=\fivei
  \textfont2=\tensy \scriptfont2=\sevensy \scriptscriptfont2=\fivesy
  \textfont3=\tenex \scriptfont3=\tenex \scriptscriptfont3=\tenex
  \def\it{\fam\itfam\tenit}%
  \textfont\itfam=\tenit
  \def\sl{\fam\slfam\tensl}%
  \textfont\slfam=\tensl
  \def\bf{\fam\bffam\tenbf}%
  \textfont\bffam=\tenbf \scriptfont\bffam=\sevenbf
   \scriptscriptfont\bffam=\fivebf
  \def\tt{\fam\ttfam\tentt}%
  \textfont\ttfam=\tentt
  \tt \ttglue=.5em plus.25em minus.15em
  \normalbaselineskip=12pt
  \def\MF{{\manual META}\-{\manual FONT}}%
  \let\sc=\eightrm
  \let\big=\tenbig
  \setbox\strutbox=\hbox{\vrule height8.5pt depth3.5pt width\z@}%
  \normalbaselines\rm}
%
\def\ninepoint{\def\rm{\fam0\ninerm}%
  \textfont0=\ninerm \scriptfont0=\sixrm \scriptscriptfont0=\fiverm
  \textfont1=\ninei \scriptfont1=\sixi \scriptscriptfont1=\fivei
  \textfont2=\ninesy \scriptfont2=\sixsy \scriptscriptfont2=\fivesy
  \textfont3=\tenex \scriptfont3=\tenex \scriptscriptfont3=\tenex
  \def\it{\fam\itfam\nineit}%
  \textfont\itfam=\nineit
  \def\sl{\fam\slfam\ninesl}%
  \textfont\slfam=\ninesl
  \def\bf{\fam\bffam\ninebf}%
  \textfont\bffam=\ninebf \scriptfont\bffam=\sixbf
   \scriptscriptfont\bffam=\fivebf
  \def\tt{\fam\ttfam\ninett}%
  \textfont\ttfam=\ninett
  \tt \ttglue=.5em plus.25em minus.15em
  \normalbaselineskip=11pt
  \def\MF{{\manual hijk}\-{\manual lmnj}}%
  \let\sc=\sevenrm
  \let\big=\ninebig
  \setbox\strutbox=\hbox{\vrule height8pt depth3pt width\z@}%
  \normalbaselines\rm}
%
\def\ttninepoint{\def\rm{\fam0\ninerm}%
  \textfont0=\ninerm \scriptfont0=\sixrm \scriptscriptfont0=\fiverm
  \textfont1=\ninei \scriptfont1=\sixi \scriptscriptfont1=\fivei
  \textfont2=\ninesy \scriptfont2=\sixsy \scriptscriptfont2=\fivesy
  \textfont3=\tenex \scriptfont3=\tenex \scriptscriptfont3=\tenex
  \def\it{\fam\itfam\nineit}%
  \textfont\itfam=\nineit
  \def\sl{\fam\slfam\ninesl}%
  \textfont\slfam=\ninesl
  \def\bf{\fam\bffam\ninebf}%
  \textfont\bffam=\ninebf \scriptfont\bffam=\sixbf
   \scriptscriptfont\bffam=\fivebf
  \def\tt{\fam\ttfam\ninett}%
  \textfont\ttfam=\ninett
  \normalbaselineskip=11pt
  \normalbaselines\rm}
%
\def\eightpoint{\def\rm{\fam0\eightrm}%
  \textfont0=\eightrm \scriptfont0=\sixrm \scriptscriptfont0=\fiverm
  \textfont1=\eighti \scriptfont1=\sixi \scriptscriptfont1=\fivei
  \textfont2=\eightsy \scriptfont2=\sixsy \scriptscriptfont2=\fivesy
  \textfont3=\tenex \scriptfont3=\tenex \scriptscriptfont3=\tenex
  \def\it{\fam\itfam\eightit}%
  \textfont\itfam=\eightit
  \def\sl{\fam\slfam\eightsl}%
  \textfont\slfam=\eightsl
  \def\bf{\fam\bffam\eightbf}%
  \textfont\bffam=\eightbf \scriptfont\bffam=\sixbf
   \scriptscriptfont\bffam=\fivebf
  \def\tt{\fam\ttfam\eighttt}%
  \textfont\ttfam=\eighttt
  \tt \ttglue=.5em plus.25em minus.15em
  \normalbaselineskip=9pt
  \def\MF{{\manual opqr}\-{\manual stuq}}%
  \let\sc=\sixrm
  \let\big=\eightbig
  \setbox\strutbox=\hbox{\vrule height7pt depth2pt width\z@}%
  \normalbaselines\rm}
%
\def\tenmath{\tenpoint\fam-1 } % use after $ in ninepoint sections
\def\tenbig#1{{\hbox{$\left#1\vbox to8.5pt{}\right.\n@space$}}}
\def\ninebig#1{{\hbox{$\textfont0=\tenrm\textfont2=\tensy
  \left#1\vbox to7.25pt{}\right.\n@space$}}}
\def\eightbig#1{{\hbox{$\textfont0=\ninerm\textfont2=\ninesy
  \left#1\vbox to6.5pt{}\right.\n@space$}}}
%
% (2) Inweb headings
% `\weavesection` and its heavier friends are Inweb sections, the grander the
% more `s` suffixes they have: untitled paragraph, titled paragraph, section, chapter.
%
\def\weavesectionsss#1#2#3#4#5{\vfill\eject\noindent\chaptertitletext{#2}\mark{#3}\vskip 0.5in\par\noindent\titlepage}
\def\weavesectionss#1#2#3#4#5{\vfill\eject\noindent\sectiontitletext{#2\hfill #5}\mark{#3}\bigskip\par\noindent\titlepage}
\def\weavesections#1#2#3#4#5{\bigskip\goodbreak\noindent{\bf #4#1. #2.}\mark{#3}\quad}
\def\weavesection#1#2#3#4#5{\bigskip\goodbreak\noindent{\bf #4#1.}\mark{#3}\quad}
%
% The `t` variants are for themed extracts.
\def\tweavesections#1#2#3#4#5{\bigskip\goodbreak\noindent{\bf #5.#4#1. #2.}\quad}
\def\tweavesection#1#2#3#4#5{\bigskip\goodbreak\noindent{\bf #5.#4#1.}\quad}
%
% The `ns` variants have less vertical spacing before them.
\def\nsweavesections#1#2#3#4#5{\noindent{\bf #4#1. #2.}\mark{#3}\quad}
\def\nsweavesection#1#2#3#4#5{\noindent{\bf #4#1.}\mark{#3}\quad}
%
% (3) Strikethrough
% Plain TeX does not support text struck through, out of the box: it turns out
% to be quite tricky to get right.
%
\catcode`@=11
\def\strikethrough#1{\leavevmode\strikethr@ugh#1 \relax}
\def\strikethr@ugh#1 {\strike@word{#1}\futurelet\next\strike@next}
\def\strike@next{\ifx\next\relax\else\strike@space\expandafter\strikethr@ugh\fi}
\def\strike@word#1{\setbox0=\hbox{#1}\rlap{\strike@rule width\wd0}\box0}
\def\strike@space{\leaders\strike@rule\sp@ce}
\def\strike@rule{\vrule height2.4pt depth-2pt}
\def\sp@ce{\hskip\fontdimen2\font plus\fontdimen3\font minus\fontdimen4\font}
\catcode`@=12
%
% (4) Nested lists, enumerated and unenumerated
% We don't want to have to rely on Plain TeX's `\item` and `\itemitem`.
% Instead, all items are `\item` (which is redefined as appropriate for the
% context) and lists begin and end with `\beginenumerate` and `\endenumerate`
% if they are to be numbered, or `\beginunenumerate` and `\endunenumerate` if not;
% and those can be fairly arbitrarily nested.
%
\catcode`@=11
\newcount\c@listlevel
\newcount\c@item
\newdimen\enumerateindent
%
\def\init@str{} % the label of every item
\def\beginenumerate{%
  \begingroup
  \let\item\enumerate@item
  \advance\c@listlevel\@ne
  \ifnum\c@listlevel=\@ne
    \vskip.5\baselineskip
  \fi
  \ifnum\c@listlevel>\@ne
    \edef\init@str{}%
  \fi
  \c@item\z@
}
%
\def\endenumerate{%
  \par\endgroup
  \ifnum\c@listlevel=\z@
    \vskip.5\baselineskip
  \fi}
%
\def\beginunenumerate{%
  \begingroup
  \let\item\unenumerate@item
  \advance\c@listlevel\@ne
  \ifnum\c@listlevel=\@ne
    \vskip.5\baselineskip
  \fi
  \c@item\z@
}
%
\def\endunenumerate{%
  \par\endgroup
  \ifnum\c@listlevel=\z@
    \vskip.5\baselineskip
  \fi}
%
\def\enumerate@item{%
  \advance\c@item\@ne
  \enumerateindent=\c@listlevel\parindent
  \ifvmode\else\par\fi
  % if an \item has more than one paragraph only the first one have to be preceded by a label;
  % so \@itemlabel redefine itself the first time is expanded.
  \def\@itemlabel{%
    \llap{\init@str\number\c@item.\quad}%
    \def\@itemlabel{}%
  }%
  \everypar{\setbox0\lastbox
    \hangindent\enumerateindent \hangafter\z@
    \@itemlabel}}
%
\def\unenumerate@item{%
  \advance\c@item\@ne
  \enumerateindent=\c@listlevel\parindent
  \ifvmode\else\par\fi
  \everypar{\setbox0\lastbox
    \hangindent\enumerateindent \hangafter\z@
    $\bullet$~}}
%
\catcode`@=12
%
% (5) Verbatim code samples
%
\chardef\other=12
\def\ttverbatim{\begingroup\ttninepoint\frenchspacing
  \catcode`\\=\other
  \catcode`\{=\other
  \catcode`\}=\other
  \catcode`\$=\other
  \catcode`\&=\other
  \catcode`\#=\other
  \catcode`\%=\other
  \catcode`\~=\other
  \catcode`\_=\other
  \catcode`\^=\other
  \obeyspaces \obeylines \tt}
%
\def\|{\leavevmode\hbox{\tt\char`\|}} % vertical line
\outer\def\begintt{$$\let\par=\endgraf \ttverbatim \parskip=\z@
  \catcode`\|=0 \rightskip-5pc \ttfinish}
{\catcode`\|=0 |catcode`|\=\other % | is temporary escape character
  |obeylines % end of line is active
  |gdef|ttfinish#1^^M#2\endtt{#1|vbox{#2}|endgroup|tenpoint$$}}
\catcode`\|=\active
%
{\obeylines \gdef|{\ttverbatim \spaceskip\ttglue \let^^M=\  \let|=\endgroup}}
%
\catcode`@=11
\def\beginlines{\par\begingroup\nobreak\medskip\parindent\z@ \obeylines\nobreak \everypar{\strut}}
\def\beginlinestighter{\par\begingroup\nobreak\parindent=0pt \kern1pt\nobreak \obeylines \everypar{\strut}}
\def\endlines{\kern1pt\endgroup\medbreak\noindent}
\catcode`@=12
%
% (6) Preform grammar document macros
% These are no longer really used, but are brief enough to be worth keeping around
% in case of an underdog comeback story.
%
\def\[#1]{\silenttrue\xref|#1|\thinspace{\tt#1}\thinspace} % keyword in syntax
\def\beginsyntax{\endgraf\nobreak\medskip
  \begingroup \catcode`<=13 \catcode`[=13
  \let\par=\endsyntaxline \obeylines}
\def\endsyntaxline{\futurelet\next\syntaxswitch}
\def\syntaxswitch{\ifx\next\<\let\next=\syntaxrule
  \else\ifx\next\endsyntax\let\next=\endgroup
  \else\let\next=\continuerule\fi\fi \next}
\def\continuerule{\hfil\break\indent\qquad}
\def\endsyntax{\medbreak\noindent}
{\catcode`<=13 \catcode`[=13
  \global\let<=\< \global\let[=\[
  \gdef\syntaxrule<#1>{\endgraf\indent\silentfalse\xref\<#1>}}
\def\is{\ $\longrightarrow$ }
\def\alt{\ $\vert$ }
\def\cast#1#2{{#1}$\;\rightarrow\;${#2}}
\def\comp#1#2{{#1}$\;\sim\;${#2}}
\def\al{$\alpha$}
\def\be{$\beta$}
\def\nonterminal#1{$\langle${\xreffont\pdfliteral direct{1 1 0 0 k}#1\special{PDF:0 g}}$\rangle$}
\def\fixed#1{{\bf #1}}
\def\from{\par\qquad\quad$\leftarrow$\quad}
\def\continuation{\par\qquad\quad\phantom{$\leftarrow$}\quad}
%
% (7) Miscellaneous math macros which seem to be expected by MathJax users,
% even though they are not in Plain TeX
%
\def\frac#1#2{{{#1}\over{#2}}}
\def\cfrac#1#2{{{#1}\over{#2}}}
\def\dfrac#1#2{{{#1}\over{#2}}}
\def\tfrac#1#2{{{#1}\over{#2}}}
\def\mathit#1{{\hbox{\it #1}}}
\def\mathrm#1{{\hbox{\rm #1}}}
\def\mathbf#1{{\hbox{\bf #1}}}
\def\text#1{{\hbox{#1}}}
\def\lvert{|}
\def\rvert{|}
\def\implies{\Rightarrow}
%
% (8) And finally, something other than macros and fonts: the page design.
% First, header and footer:
%
\nopagenumbers
\newif\iftitle
\def\titlepage{\global\titletrue}
\headline={\iftitle\hfil\global\titlefalse\else\tenrm\hfil{\sourcefont \firstmark\qquad\the\pageno}\fi}
\titlepage
%
% Tolerances for spacing:
%
\raggedbottom
\tolerance=20000
%
% End of `inweb-macros.tex`
